% s-049-p049-m002; original PDF 49, 50
\Aphorism{s-049-p049-m002}{40}
Damit es sich in der „Wahrheit“ leben liesse, damit wir auch nur das Recht auf diese \Sperr{Bequemlichkeit} ansprechen dürften, müsste der Mensch im letzten Inneren anders beschaffen sein als er beschaffen ist: weniger Laune und mehr Wille, weniger Nerven und mehr Knochen. Wo wir uns einstellen, wird uns zuerst Dauer, Gleichgewicht, Berechenbarkeit abverlangt, bis wir sie schliesslich selber von uns verlangen; wir billigen den ästhetischen Stimmungsmenschen nicht, der sich durch eine zerfahrene Praxis darzustellen sucht — und doch, was liegt uns im Grunde näher als solche Stimmungsmenschen zu sein und Berufe, Aufenthalte, Ideale leichten Herzens zu wechseln? Hier tritt die Lüge, die \Foreign{pia fraus} in ihre Rechte, als einzige Möglichkeit, unter ausgelebten Verhältnissen, in zerrissenem Zusammenhange weiter zu existiren; und wo nur irgend Bestand und Beständigkeit verehrt wird, sollte auch der heiligen Lüge ein Altar geweiht sein. Es ist unmöglich, die Dauer einer Fähigkeit, das Leben einer Überzeugung, die Treue einer Neigung zu verbürgen; aber es ist stets möglich, den Schein dieser Dinge im Zusammenleben mit Anderen zu wahren. Wir können für unsere Oberfläche einstehen, nicht für unser Inneres; wir können Wirkungen versprechen, nicht auch Ursachen. Man sollte eine Rangordnung der Menschen unter dem Gesichtspunkte aufstellen, worin und zu welchen Zwecken sie heucheln; gewiss würde sie die edelsten wie die niedrigsten Typen umfassen. Die „wahrhaftigen“ Menschen würden von diesem Wettbewerb auszuschliessen sein, als Schosskinder des Glücks, denen das Leben Schwierigkeiten und Conflicte erspart hat, oder als Trunkenbolde des Ideals, die sich nicht wehethun, wenn sie fallen, aber auch nichts davon merken, dass sie Andere anstossen. Dem Wahrhaftigen fehlen die sympathischen Functionen, wahrscheinlich weil er nicht nöthig hatte, sie auszubilden; wer einmal Arzt, Erzieher, Liebender, Freund gewesen ist, hat auch seine Wirkung auf Andere bemessen, d. h. hat lügen gelernt.
\par

